\documentclass[12pt]{amsart}
%prepared in AMSLaTeX, under LaTeX2e
\addtolength{\oddsidemargin}{-.65in} 
\addtolength{\evensidemargin}{-.65in}
\addtolength{\topmargin}{-.55in}
\addtolength{\textwidth}{1.3in}
\addtolength{\textheight}{.9in}

\renewcommand{\baselinestretch}{1.05}

\usepackage{verbatim} % for "comment" environment

\usepackage{palatino}

\newtheorem*{thm}{Theorem}
\newtheorem*{defn}{Definition}
\newtheorem*{example}{Example}
\newtheorem*{problem}{Problem}
\newtheorem*{remark}{Remark}

\usepackage{fancyvrb,xspace,dsfont}

\usepackage[final]{graphicx}

% macros
\usepackage{amssymb}
\usepackage[dvipsnames]{xcolor}

\usepackage{hyperref}
\hypersetup{pdfauthor={Ed Bueler},
            pdfcreator={pdflatex},
            colorlinks=true,
            citecolor=blue,
            linkcolor=red,
            urlcolor=blue,
            }

\newcommand{\bn}{\mathbf{n}}
\newcommand{\br}{\mathbf{r}}
\newcommand{\bv}{\mathbf{v}}
\newcommand{\bx}{\mathbf{x}}
\newcommand{\by}{\mathbf{y}}

\newcommand{\bJ}{\mathbf{J}}
\newcommand{\bV}{\mathbf{V}}

\newcommand{\cB}{\mathcal{B}}
\newcommand{\cD}{\mathcal{D}}
\newcommand{\cE}{\mathcal{E}}
\newcommand{\cF}{\mathcal{F}}
\newcommand{\cH}{\mathcal{H}}
\newcommand{\cL}{\mathcal{L}}
\newcommand{\cM}{\mathcal{M}}
\newcommand{\cR}{\mathcal{R}}
\newcommand{\cV}{\mathcal{V}}
\newcommand{\cW}{\mathcal{W}}

\newcommand{\CC}{\mathbb{C}}
\newcommand{\NN}{\mathbb{N}}
\newcommand{\RR}{\mathbb{R}}
\newcommand{\ZZ}{\mathbb{Z}}

\renewcommand{\Im}{\mathrm{Im}}
\renewcommand{\Re}{\mathrm{Re}}

\newcommand{\eps}{\epsilon}
\newcommand{\lam}{\lambda}
\newcommand{\lap}{\triangle}

\newcommand{\Div}{\nabla\cdot}
\newcommand{\grad}{\nabla}

\newcommand{\one}{\mathds{1}}

\newcommand{\ip}[2]{\ensuremath{\left<#1,#2\right>}}

\newcommand{\image}{\operatorname{im}}
\newcommand{\onull}{\operatorname{null}}
\newcommand{\rank}{\operatorname{rank}}
\newcommand{\range}{\operatorname{range}}
\newcommand{\trace}{\operatorname{tr}}
\newcommand{\Span}{\operatorname{span}}

\newcommand{\prob}[1]{\bigskip\noindent\textbf{#1.}\quad }

\newcommand{\pts}[1]{(\emph{#1 pts}) }
\newcommand{\epart}[1]{\medskip\noindent\textbf{(#1)}\quad }
\newcommand{\ppart}[1]{\,\textbf{(#1)}\quad }

\newcommand{\ds}{\displaystyle}

\newcommand{\nex}{\medskip\noindent}


\begin{document}
\scriptsize \noindent Math 617 Functional Analysis (Bueler) \hfill \emph{assigned 25 March 2026}
\normalsize\medskip

\Large\centerline{\textbf{Assignment 7}}
\large
\medskip

\centerline{\textbf{Due Friday 10 April at the beginning of class}}
\medskip
\normalsize

\thispagestyle{empty}

\bigskip
\noindent This Assignment is based on the slides, especially \href{https://bueler.github.io/fa/assets/slides/S26/week10.pdf}{weeks 10 \& 11} and \href{https://bueler.github.io/fa/assets/slides/S26/week12.pdf}{week 12}.

\newcommand{\augment}[1]{\quad\,\strut \begin{minipage}[t]{0.7\textwidth} \emph{#1}\end{minipage}}

\newcommand{\augpart}[1]{\emph{\textbf{(#1)}}}

\medskip
\noindent \textsc{Do the following problems.}


\prob{P17} \emph{Compare the ``main idea'' about the definition of $C_c^\infty(\Omega)$ in the weeks 10 \& 11 slides.}

\epart{a} Suppose $\Omega\subset \RR^d$ is open.  Suppose $g\in L_{loc}^1(\Omega)$ and $\varphi \in C_c^\infty(\Omega)$.  Prove that $g D^\alpha\varphi \in L^1(\Omega)$, thus that $\int_\Omega g D^\alpha \varphi\,dm$ is well-defined.

\epart{b}  Show that $C(\Omega)\subset L_{loc}^1(\Omega)$.  Then, for each dimension $d$, construct an example showing that ``$C(\Omega)\subset L^1(\Omega)$'' is false.

\prob{P18} \emph{Here is a practical case where the weak derivative concept is appropriate in 1D.}

\medskip \noindent Suppose $a=t_0<t_1<\dots<t_n=b$ is a mesh on a closed interval $[a,b]\subset \RR^1$, and consider values $y_0,y_1,\dots,y_n\in\RR^1$.  An interpolating \emph{cubic spline} is a function $s\in C^1([a,b])$ such that $s(x)$ on each subinterval is a cubic polynomial,
    $$s(x)=a_k+b_k (x-x_{k-1}) + c_k (x-x_{k-1})^2 + d_k (x-x_{k-1})^3 \quad \text{ on } [t_{k-1},t_k],$$
for $k=1,\dots,n$, and so that $s(t_k) = y_k$, for $k=0,1,\dots,n$.  It is a theorem, which you do \emph{not} need to prove, that for the data $(t_k,y_k)$ there is a unique function $s(x)$ with these properties.  Give a formula for the weak second derivative $s''$, on $[a,b]$, and \underline{show} that it satisfies the definition of a weak derivative.  Approximately sketch the spline $s(x)$ and its second derivative $s''(x)$ in a simple case, for instance with $n=3$. (\emph{Note the key fact that $s$ and $s'$ are \emph{continuous} on $[a,b]$.  Generally $s''$ is not continuous, or even defined at all points in $[a,b]$.}) 

\prob{P19}  Recall that, if $f^{(\ell)}$ denotes the weak derivative of order $\ell$, then the usual definition of the norm in $H^k(-\pi,\pi)$ is \quad  $\ds (\|f\|_{H^k})^2 = \sum_{\ell=0}^k \int_{-\pi}^\pi |f^{(\ell)}(x)|^2\,dx$.

\epart{a}  Suppose $f\in H^1(-\pi,\pi)$.  Using the formulas for Fourier series,\footnote{See ``recall Fourier series in $L^2(-\pi,\pi)$'' in the \href{https://bueler.github.io/fa/assets/slides/S26/week12.pdf}{week 12} slides.} prove the derivative formula
	$$f'(x) = \frac{1}{\sqrt{2\pi}} \sum_{n\in\ZZ} i n c_n e^{inx},$$
with convergence in $L^2$, where $c_n$ are the Fourier coefficients of $f$.

\epart{b}  Show that
	$$\|f\|_{H^k}^2 = \sum_{n\in\ZZ} \left(\sum_{j=0}^k n^{2j}\right) |c_n|^2.$$
(\emph{There is no need to justify derivative formulas, beyond what you did in part} \textbf{(a)}.)

\prob{P20}  FIXME exactly compute $L^2(D)$ and $H^1(D)$ norms of $g(x)=\ln(r)$ and $p(x)=\ln(|\ln(r)|)$; plot $p(r)$ on $(0,1)$

\prob{P21}  FIXME shows that $H_0^1(\Omega)$ is a proper closed subspace if $\Omega=(a,b)$: $f(x)=1 \in H^1(a,b)$ but $f \notin H_0^1(a,b)$ because of the FTC

\prob{P22}  FIXME $p(z)=z^3 - 5 z^2 + 2 z - 1$  show real and imaginary parts are harmonic

\end{document}
